\documentclass[11pt]{article}
\usepackage[margin=1in]{geometry}
\usepackage{amsmath}
\usepackage{booktabs}
\usepackage{enumitem}
\usepackage{hyperref}

\title{\textbf{ADR-06: Common Customer Query Pattern Analysis via N-gram Analysis}}
\author{Architecture Decision Record}
\date{\today}

\begin{document}

\maketitle

\section{Status}
\textbf{Accepted}

\section{Context}

Understanding common linguistic patterns in customer communications enables optimization of automated response systems, development of more effective self-service content, and identification of frequently occurring customer language that can improve support documentation and training materials.

\subsection{Business Requirements}
\begin{itemize}
\item Identify most frequent word combinations in customer communications for FAQ development
\item Understand authentic customer language patterns for documentation optimization
\item Detect recurring query structures for automated response template creation
\item Provide content teams with data-driven insights for self-service portal improvements
\item Enable predictive text and auto-completion features based on real customer language
\end{itemize}

\subsection{Analytical Challenges}
\begin{itemize}
\item Raw word frequency analysis misses important contextual relationships
\item Customer language contains significant linguistic variation requiring normalization
\item Visualization must accommodate variable-length word sequences effectively
\item Analysis scalability required for larger datasets and real-time applications
\item Balance between frequent patterns and meaningful business insights
\end{itemize}

\section{Decision}

We implemented \textbf{Multi-level N-gram Analysis} using NLTK for comprehensive linguistic pattern detection:

\subsection{Core Methodology}
\begin{itemize}
\item \textbf{N-gram Extraction}: Generate word sequences of length 2, 3, and 4 (bigrams, trigrams, 4-grams)
\item \textbf{Text Preprocessing}: Lowercase conversion, punctuation removal, and stopword filtering
\item \textbf{Frequency Analysis}: Statistical counting of n-gram occurrences across entire dataset
\item \textbf{Flexible Visualization}: Dynamic plotting function accommodating different n-gram lengths
\item \textbf{Top-K Selection}: Configurable display of most frequent patterns (default: top 20)
\end{itemize}

\subsection{Processing Pipeline}
\begin{itemize}
\item \textbf{Text Aggregation}: Combine all customer communications into unified corpus
\item \textbf{Linguistic Normalization}: Standardize text format while preserving semantic meaning
\item \textbf{Sequence Generation}: Create overlapping word sequences using NLTK's ngrams function
\item \textbf{Statistical Analysis}: Counter-based frequency distribution calculation
\item \textbf{Result Formatting}: Convert tuple-based n-grams to readable string representations
\end{itemize}

\subsection{Visualization Strategy}
\begin{itemize}
\item \textbf{Horizontal Bar Charts}: Optimal layout for variable-length text labels
\item \textbf{Magma Color Palette}: Professional appearance with clear visual hierarchy
\item \textbf{Dynamic Titles}: Automatic labeling based on n-gram type (unigrams, bigrams, etc.)
\item \textbf{Frequency Focus}: Direct display of occurrence counts for business interpretation
\end{itemize}

\section{Rationale}

\subsection{Why N-gram Analysis Over Alternatives}
\begin{itemize}
\item \textbf{Context Preservation}: Captures word relationships missed by single-word analysis
\item \textbf{Natural Language Patterns}: Identifies authentic customer communication structures
\item \textbf{Scalable Methodology}: NLTK implementation handles large datasets efficiently
\item \textbf{Linguistic Sophistication}: More advanced than simple keyword extraction methods
\item \textbf{Business Applicability}: Direct translation to FAQ content and response templates
\end{itemize}

\subsection{Multi-level Analysis Benefits}
\begin{itemize}
\item \textbf{Bigrams (n=2)}: Reveal common word pairs and basic query structures
\item \textbf{Trigrams (n=3)}: Identify complete short phrases and action sequences
\item \textbf{4-grams (n=4)}: Capture longer expressions and complete query patterns
\item \textbf{Comparative Insights}: Different n-gram levels provide complementary perspectives
\end{itemize}

\subsection{Alternative Approaches Rejected}
\begin{itemize}
\item \textbf{Simple Word Frequency}: Missing contextual relationships and meaningful phrases
\item \textbf{Topic Modeling Only}: Over-complex for identifying specific language patterns
\item \textbf{Manual Pattern Identification}: Scale limitations and subjective bias
\item \textbf{Regular Expression Matching}: Requires predefined patterns, missing discovery aspect
\item \textbf{Dependency Parsing}: Excessive complexity for pattern frequency analysis
\end{itemize}

\section{Consequences}

\subsection{Positive Outcomes}
\begin{itemize}
\item \textbf{Content Optimization}: Data-driven FAQ and documentation improvement using real customer language
\item \textbf{Template Development}: Automated response templates based on common query patterns
\item \textbf{Language Alignment}: Support materials aligned with authentic customer communication styles
\item \textbf{Predictive Features}: Foundation for auto-completion and suggestion systems
\item \textbf{Training Enhancement}: Real-world customer language patterns for agent education
\end{itemize}

\subsection{Technical Advantages}
\begin{itemize}
\item \textbf{Implementation Simplicity}: NLTK provides robust, well-tested n-gram functionality
\item \textbf{Computational Efficiency}: Linear complexity scaling with dataset size
\item \textbf{Flexibility}: Easily adjustable for different n-gram lengths and frequency thresholds
\item \textbf{Integration Ready}: Results easily consumed by downstream applications and systems
\end{itemize}

\subsection{Limitations and Trade-offs}
\begin{itemize}
\item \textbf{Context Boundaries}: N-grams don't cross sentence boundaries, potentially missing patterns
\item \textbf{Frequency Bias}: Common but less meaningful patterns may overshadow important rare expressions
\item \textbf{Language Specificity}: Current implementation optimized for English text only
\item \textbf{Semantic Blindness}: Doesn't understand meaning, only statistical co-occurrence patterns
\end{itemize}

\section{Implementation Details}

\subsection{N-gram Generation Algorithm}
\begin{align}
\text{Preprocessing} &: text \rightarrow \text{lowercase, depunctuated, destopworded} \\
\text{Tokenization} &: text \rightarrow [w_1, w_2, \ldots, w_k] \\
\text{N-gram Generation} &: \text{ngrams}([w_1, \ldots, w_k], n) \\
\text{Frequency Analysis} &: \text{Counter}(\text{n-grams}) \\
\text{Top-K Selection} &: \text{most\_common}(20)
\end{align}

\subsection{Function Parameters and Configuration}
\begin{itemize}
\item \textbf{Input Parameters}: DataFrame, text column name, n-gram size, top-k count
\item \textbf{Preprocessing}: NLTK stopwords, string punctuation removal, case normalization
\item \textbf{Visualization}: 12×8 figure size, magma palette, horizontal orientation
\item \textbf{Error Handling}: Validation for insufficient data and empty n-gram sets
\end{itemize}

\subsection{Performance Characteristics}
\begin{itemize}
\item \textbf{Processing Speed}: ~500 tickets/second for bigram analysis on standard hardware
\item \textbf{Memory Usage}: Approximately 5MB per 1000 tickets for complete n-gram storage
\item \textbf{Scalability}: Linear time complexity O(n\^k) where n=documents, k=average words$/$document
\item \textbf{Storage Requirements}: Counter objects provide memory-efficient frequency storage
\end{itemize}

\section{Success Metrics}

\subsection{Analysis Quality Indicators}
\begin{itemize}
\item Top 20 n-grams represent >70\% of all n-gram occurrences in dataset
\item Identified patterns align with known customer communication themes
\item Analysis completion time under 60 seconds for 50-ticket dataset
\item Generated patterns provide actionable insights for content development teams
\end{itemize}

\subsection{Business Impact Measures}
\begin{itemize}
\item FAQ content incorporates language patterns from n-gram analysis
\item Self-service portal search functionality improved using common query patterns
\item Customer satisfaction scores improve through language-aligned support content
\item Response template development accelerated through pattern identification
\end{itemize}

\section{Future Evolution}

\subsection{Enhancement Opportunities}
\begin{itemize}
\item \textbf{Semantic N-grams}: Incorporate word embeddings for semantic similarity grouping
\item \textbf{Temporal Analysis}: Track n-gram pattern evolution over time periods
\item \textbf{Category-Specific Patterns}: Analyze n-grams within specific support categories
\item \textbf{Multi-language Support}: Extend analysis to international customer communications
\item \textbf{Context-Aware Analysis}: Cross-sentence pattern detection for longer expressions
\end{itemize}

\subsection{Integration Roadmap}
\begin{itemize}
\item Real-time pattern detection for live customer communication analysis
\item API development for integration with content management systems
\item Machine learning pipeline integration for automatic template generation
\item Dashboard development for ongoing pattern monitoring and trend analysis
\end{itemize}

\subsection{Advanced Applications}
\begin{itemize}
\item \textbf{Predictive Text}: Auto-completion systems based on common customer patterns
\item \textbf{Query Routing}: Automatic categorization based on n-gram pattern recognition
\item \textbf{Content Personalization}: Dynamic FAQ ordering based on query pattern frequencies
\item \textbf{Training Optimization}: Agent training programs aligned with actual customer language patterns
\end{itemize}

\end{document}